\documentclass[11pt]{amsart}
\usepackage[T1]{fontenc}
\usepackage[utf8]{inputenc}
\usepackage{lmodern}
\usepackage{amsmath,amssymb,mathtools}
\usepackage{microtype}
\usepackage[a4paper,margin=32mm]{geometry}
\usepackage[hidelinks,hypertexnames=false]{hyperref}
\newtheorem{theorem}{Theorem}[section]
\newtheorem{proposition}[theorem]{Proposition}
\newtheorem{lemma}[theorem]{Lemma}
\newtheorem{corollary}[theorem]{Corollary}
\theoremstyle{remark}\newtheorem{remark}[theorem]{Remark}
\numberwithin{equation}{section}
\title[Steklov inequalities for quadrilaterals]{A sharp Steklov inequality for quadrilaterals\protect\\ by interpolation along a diagonal}
\author{Zhuo Cheng}
\author{Yi'an Zhang}
\date{}
\keywords{Steklov eigenvalue, quadrilateral, quantitative stability, spectral optimization}
\begin{document}
\begin{abstract}
Let $\tau\in(0,\pi/2)$ be the unique solution of $\tan\tau\tanh\tau=1$. We prove that every simple nondegenerate planar quadrilateral $Q$ of perimeter $P(Q)$ satisfies $P(Q)\sigma_1(Q)\le 8\tau\tanh\tau$, where $\sigma_1(Q)$ denotes the first nonzero Steklov eigenvalue, with equality if and only if $Q$ is a square. For convex quadrilaterals with four genuine corners, we establish a quantitative stability estimate that is quadratic in the displacement of the diagonal midpoints and linear in the affine deformation of the associated parallelogram. In the convex case, the proof employs a pair of trial functions interpolating between the first square eigenfunctions and the coordinate map via a two-triangle parametrization along a diagonal; this yields an exact completion of the Dirichlet energy, after which a boundary covariance estimate and a square completion control the full displacement range. Concave quadrilaterals satisfy $P(Q)\sigma_1(Q)<3\sqrt{3}$.
\end{abstract}
\maketitle
\raggedbottom

\section{Introduction}

For a bounded Lipschitz domain $\Omega\subset\mathbb R^2$ of perimeter $P(\Omega)$, the first nonzero Steklov eigenvalue admits the variational characterization
\[
\sigma_1(\Omega)
=\inf_{\substack{w\in H^1(\Omega)\\ \mathcal M_\Omega(w)>0}}
\frac{\mathcal E_\Omega(w)}{\mathcal M_\Omega(w)},
\]
where the Dirichlet energy and boundary variance are given by
\[
\mathcal E_\Omega(w)=\int_\Omega|\nabla w|^2,
\qquad
\mathcal M_\Omega(w)=\int_{\partial\Omega}w^2\,ds
-\frac1{P(\Omega)}\left(\int_{\partial\Omega}w\,ds\right)^2,
\]
with $ds$ denoting arclength measure on $\partial\Omega$. Because the perimeter-normalized product $\mathcal J(\Omega)=P(\Omega)\sigma_1(\Omega)$ is invariant under similarities, it is the natural scale-free functional to maximize at fixed perimeter.

For simply connected planar domains, Weinstock~\cite{Weinstock1954}
proved $\mathcal J(\Omega)\le2\pi$, with equality precisely for disks.
Hersch, Payne, and Schiffer~\cite{HerschPayneSchiffer1974} obtained
inequalities for products of higher Steklov eigenvalues. Girouard and
Polterovich~\cite{GirouardPolterovich2010} proved sharpness for individual
eigenvalues and products of consecutive ones through domains degenerating
into equal disks. Under a volume constraint, Brock's
inequality~\cite{Brock2001} identifies the ball as the maximizer of the
first nonzero eigenvalue. Bucur, Ferone, Nitsch, and
Trombetti~\cite{BucurEtAl2021} proved the corresponding result under a
boundary-measure constraint in the convex class.

The geometric theory also connects extremal Steklov metrics with free
boundary minimal surfaces, through the work of Fraser and
Schoen~\cite{FraserSchoen2011,FraserSchoen2016}.
Continuity of variational eigenvalues and homogenization led Girouard,
Karpukhin, and Lagac\'e~\cite{GirouardKarpukhinLagace2021} to sharp bounds
for planar domains with unrestricted boundary topology. The surveys
\cite{GirouardPolterovich2017,ColboisEtAl2024} describe these developments
and their place in spectral geometry.

Quantitative inequalities relate the spectral deficit to geometric
deformation. Brasco, De Philippis, and Ruffini~\cite{BrascoEtAl2012}
proved sharp stability for the volume-constrained inequality.
Bucur and Nahon~\cite{BucurNahon2021} identified the role of boundary
oscillations in the stability of Weinstock's perimeter inequality.

A fixed number of straight sides leads to extremal problems for regular
polygons. Their Dirichlet and torsional counterparts originate in
P\'olya's symmetrization theory~\cite{Polya1948}.
Bucur and Fragal\`a~\cite{BucurFragala2021} obtained symmetry results from
the first variations of polygonal energies, while Bogosel and
Bucur~\cite{BogoselBucur2024} reduced the polygonal Faber--Krahn problem
to finitely many certified computations. For the Steklov problem, Cheng,
Gui, Hu, Li, and Yao~\cite{ChengEtAl2026} proved that the first eigenvalue
of a regular $N$-gon of fixed perimeter increases strictly with $N$ for
every $N\ge3$, and that its first eigenspace has dimension two. The
polygonal extremal conjecture recorded in
\cite[Conjecture~1.1]{ChengEtAl2026} asks whether the regular $N$-gon
maximizes this eigenvalue among convex polygons with at most $N$ sides
and prescribed perimeter.

Recent torsion results give further comparisons within polygonal classes.
Gui, Hu, Li, and Zhang~\cite{GuiHuLiZhang2026} studied mixed torsion on
right triangles and proved strict monotonicity of Dirichlet torsional
rigidity along equal-area regular polygons. For tangential polygons, Gui,
Hu, and Li~\cite{GuiHuLi2026} established a quantitative
P\'olya--Szeg\H{o} inequality with an explicit perimeter deficit.
Gui, Hu, Li, and Yang~\cite{GuiHuLiYang2026} extended this comparison to
Robin torsional rigidity. These results quantify the effect of polygonal
geometry on variational energies.

The Steklov spectrum records the corners as well as the side lengths.
Levitin, Parnovski, Polterovich, and Sher~\cite{LevitinEtAl2022} derived
asymptotics for curvilinear polygons, and Krymski, Levitin, Parnovski,
Polterovich, and Sher~\cite{KrymskiEtAl2021} used these asymptotics in the
inverse problem. Dryden, Gordon, Moreno, Rowlett, and
Villegas-Blas established spectral finiteness for generic convex
polygons~\cite{DrydenEtAl2025} and studied spectral determination for
triangles, quadrilaterals, and regular polygons~\cite{DrydenEtAl2026}.

For optimization over variable domains, Bogosel~\cite{Bogosel2017}
proved continuity and existence results under geometric control, and
Bogosel, Bucur, and Giacomini~\cite{BogoselBucurGiacomini2017} established
existence of volume-constrained maximizers in a relaxed class.
Numerical methods for boundary eigenvalue problems and extremal Steklov
shapes were developed in~\cite{Bogosel2016,AkhmetgaliyevEtAl2017}.
On rectangles, explicit Steklov eigenfunctions also give boundary
integral representations of harmonic functions~\cite{AuchmutyCho2015}.

Let $\tau\in(0,\pi/2)$ be the unique root of $\tan\tau\tanh\tau=1$, and set $r=\tanh\tau$ and $s=\tau r=\tau\tanh\tau$. Separation of variables on the reference square $\Omega_\square=(-1,1)^2$ shows that $\sigma_1(\Omega_\square)=s$, with a two-dimensional first eigenspace \cite{Tan2017,GirouardEtAl2019}. Since $P(\Omega_\square)=8$, its normalized eigenvalue is $\mathcal J(\Omega_\square)=8s$. The sharp quadrilateral inequality is the following.

\begin{theorem}\label{thm-main}
Every simple nondegenerate quadrilateral $Q\subset\mathbb R^2$ satisfies
\[
P(Q)\sigma_1(Q)\le 8s,
\]
with equality if and only if $Q$ is a square. Consequently, up to rigid motions, the square uniquely maximizes $\sigma_1$ among simple quadrilaterals of prescribed perimeter.
\end{theorem}

Here a \emph{simple nondegenerate quadrilateral} means the polygonal domain bounded by a simple closed chain of four segments of positive length enclosing positive area; a redundant collinear vertex (in which case $Q$ is geometrically a triangle) is permitted. For convex quadrilaterals with four genuine corners, Lemma~\ref{lem-normal-form} introduces normalized midpoint coordinates $(b,\gamma,h,k)$ in which $(b,\gamma)$ encodes the aspect ratio and shear of the associated parallelogram (with $b=1$ and $\gamma=0$ when this parallelogram is a square), while $(h,k)$ measures the displacement of the diagonal midpoints. In these coordinates we obtain the following quantitative stability estimate.

\begin{samepage}
\begin{theorem}\label{thm-stability}
There exists a universal constant $c>0$ such that every convex quadrilateral $Q\subset\mathbb R^2$ with four genuine corners, expressed in the normalized midpoint coordinates $(b,\gamma,h,k)$ of Lemma~\ref{lem-normal-form}, satisfies
\[
8s-\mathcal J(Q)\ge c\min\left\{1,\,
h^2+k^2+\sqrt{(1-b)^2+\gamma^2}\right\}.
\]
\end{theorem}
\end{samepage}

Transplantation and coordinate tests are classical tools in isoperimetric
spectral inequalities~\cite{PolyaSchiffer1954,Szego1954,Weinberger1956}.

For convex quadrilaterals, an elementary diagonal test restricts the problem to a fixed range of affine parameters. In this regime, we parametrize $Q$ from $\Omega_\square$ by a bi-Lipschitz map that is affine on each of the two triangles separated by a diagonal. After factoring out the affine part of the associated parallelogram, the displacement of the diagonal midpoints enters through the scalar profile $q(x,y)=1-|x-y|$ on $\Omega_\square$, which agrees with $xy$ on $\partial\Omega_\square$, and the two triangular pieces have Jacobian determinants proportional to $1\pm(h-k)$. Pairing opposite points under central reflection then yields an exact completion of the Dirichlet energy with denominator $1-(h-k)^2$. We construct a pair of trial functions interpolating between the first square eigenfunctions and the linear coordinates with weight $\lambda=1-(h-k)^2$, which cancels this denominator.

Transporting this trial pair to physical coordinates decomposes the total Rayleigh deficit into an affine contribution carrying positive quadratic terms in $1-b$ and $\gamma$, a boundary displacement gain controlled uniformly across the full convex displacement range by a boundary variance estimate, and a mixed energy term coupling $1-b$ with the diagonal displacement. A square completion gives a nonnegative lower bound that vanishes only for the square, establishing both the sharp convex inequality and quadratic stability in $(h,k)$. The linear affine rate in Theorem~\ref{thm-stability} follows from the largest eigenvalue of the deficit form on the two-dimensional first eigenspace of the square. Its first-order part detects both $1-b$ and $\gamma$, while the displacement changes this form by $O(h^2+k^2)$. Finally, broken-line boundary variance estimates for the coordinate functions give $\mathcal J(Q)<3\sqrt{3}$ for concave quadrilaterals and $\mathcal J(Q)\le3\sqrt{3}$ for triangles with a redundant vertex.

Section~2 introduces the midpoint coordinates, the square moments and the
diagonal exclusion estimate. Section~\ref{sec-global-interpolation}
develops the diagonal interpolation and the boundary covariance estimate.
Section~4 proves the convex deficit, and Section~5 derives quantitative
stability. Section~6 treats concave quadrilaterals and triangles with a
redundant vertex. The square moment bounds are proved in
Appendix~\ref{app-square}.

\section{Geometry and square data}

For an admissible trial function define its normalized deficit by
\[
D_Q(w)=\frac{\mathcal M_Q(w)}{P(Q)}-
\frac{\mathcal E_Q(w)}{8s}.
\]
We use the same notation with two arguments for the associated symmetric
bilinear forms. The variational principle gives
\begin{equation}\label{eq-rayleigh-deficit}
(8s-\mathcal J(Q))\frac{\mathcal M_Q(w)}{P(Q)}
\ge8s\,D_Q(w).
\end{equation}
Summing this inequality over a finite collection of trials shows that
a positive total deficit implies $\mathcal J(Q)<8s$.

\begin{samepage}
\begin{lemma}[Midpoint coordinates]\label{lem-normal-form}
After a similarity and relabeling, every convex quadrilateral with four
genuine corners has cyclic vertices
\[
-e-f+g,\quad e-f-g,\quad e+f+g,\quad -e+f-g,
\]
where
\[
e=(1,0),\quad f=b(\gamma,\sqrt{1-\gamma^2}),\quad g=he+kf,
\]
\[
0<b\le1,\quad |\gamma|<1,\quad h,k\ge0,\quad h+k<1.
\]
The area and the larger diagonal length $L$ are
\[
|Q|=4b\sqrt{1-\gamma^2},\qquad
L^2=4(1+b^2+2b|\gamma|).
\]
\end{lemma}
\end{samepage}
\begin{proof}
Translate the vertices $v_1,\ldots,v_4$ so that their sum is zero, and set
\[
e=\frac{v_2+v_3-v_1-v_4}{4},\quad
f=\frac{v_3+v_4-v_1-v_2}{4},\quad
g=\frac{v_1+v_3-v_2-v_4}{4}.
\]
The shoelace identity gives signed area $4\det(e,f)$, so $e,f$ are
independent. Write $g=he+kf$. The four consecutive edge determinants,
after division by $4\det(e,f)$, are
$1-h+k$, $1+h+k$, $1+h-k$, and $1-h-k$.
Thus strict convexity is equivalent to $|h|+|k|<1$.
Reversing the order can interchange $e,f$, so arrange $|e|\ge|f|$.
The relabelings $(v_2,v_1,v_4,v_3)$ and $(v_4,v_3,v_2,v_1)$
change the signs of $k$ and $h$, respectively, without changing this
ordering. A similarity and a possible reflection give the stated form.
The side midpoints are $-f,e,f,-e$, and the diagonal midpoints are $g,-g$.
Thus $A\Omega_\square$, with $A=[e\ f]$, has the same side midpoints as $Q$.
The diagonal vectors are $2(e+f)$ and $2(f-e)$, which give the remaining formulas.
\end{proof}

We use the first square eigenfunctions
\[
u(x,y)=\frac{\sin(\tau x)\cosh(\tau y)}{\cosh\tau},\qquad
v(x,y)=u(y,x),\qquad (x,y)\in\Omega_\square.
\]
Define their boundary and energy moments by
\[
a_\square=\int_{-1}^1u(t,1)^2\,dt,\quad
b_\square=\int_{-1}^1u(1,t)^2\,dt,\quad
\ell_0=\int_{-1}^1u(1,t)\,dt,\quad
\ell_1=\int_{-1}^1t\,u(t,1)\,dt,
\]
\[
E=\int_{\Omega_\square}|\nabla u|^2,\quad
B=\int_{\Omega_\square}(\partial_yu)^2,\quad
C=\int_{\Omega_\square}
(\partial_xu\,\partial_yv+\partial_yu\,\partial_xv),
\]
\[
D_0=E-C,\qquad C_*=E-4\ell_0+4.
\]

\begin{lemma}[Square moments]\label{lem-square-model-bounds}
The square eigenvalue is $s$, and the moments above satisfy
\[
\frac{11}{16}<s<\frac7{10},\quad
\frac{49}{100}<a_\square\le b_\square<\frac{81}{100},\quad
\ell_0>\frac54,\quad \ell_0^2<\frac85,\quad
\frac9{16}<\ell_1<\frac23,
\]
\[
E<\frac95,\qquad 0<B<\frac3{20},\qquad
0<D_0<\frac18,\qquad C_*<\frac34,\qquad C>\frac85.
\]
\end{lemma}
\begin{proof}
The eigenvalue ordering follows from the separated rectangular spectrum
\cite{Tan2017,GirouardEtAl2019}. Appendix~\ref{app-square} derives the moment identities and proves the
scalar bounds.
\end{proof}

\begin{lemma}[A broken line]\label{lem-broken-line}
Let a two-segment chain join $(-1,0)$ to $(1,0)$ through $(z,y)$, and let
its length be $2p$, where $p\ge1$. For normalized arclength on that chain,
\[
\mathbb E(x,y)=\left(\frac z2(1-p^{-2}),\frac y2\right),
\qquad
\operatorname{Var}(x)=\frac13+
\frac{(p^2-1)(p^2+3)z^2}{12p^4},
\]
\[
\operatorname{Var}(y)=\frac{y^2}{12},\qquad
\operatorname{Var}(x)+\operatorname{Var}(y)
=\frac{p^2+3}{12}+\frac{(p^2-1)z^2}{4p^4}.
\]
\end{lemma}
\begin{proof}
The two segment lengths are $p+z/p$ and $p-z/p$, and
$y^2=(p^2-1)(1-z^2/p^2)$. On each segment, integrate the affine coordinates using the endpoint
average for the first moment and
$(\alpha^2+\alpha\beta+\beta^2)/3$ for the second moment of a coordinate
with endpoint values $\alpha,\beta$. Subtracting the squares of the means
gives the formulas.
\end{proof}

\begin{lemma}[Diagonal exclusion]\label{lem-diagonal-exclusion}
For either diagonal of length $L$ in a simple quadrilateral,
\[
\mathcal J(Q)\le\frac{12|Q|}{L^2}.
\]
In the convex normal form, if $b\le13/20$ or
$|\gamma|\ge1/11$, then $\mathcal J(Q)\le11/2<8s$.
\end{lemma}
\begin{proof}
The two boundary arcs joining the diagonal endpoints are broken lines.
By Lemma~\ref{lem-broken-line}, each has variance at least $L^2/12$
in the diagonal direction. The variance of the boundary measure is the weighted average of these
variances plus a nonnegative term from the difference of the two means.
The linear coordinate in that direction has energy $|Q|$, proving the
first estimate. With $\rho=|\gamma|$, the normal form gives
\[
\mathcal J(Q)\le
\frac{12b\sqrt{1-\rho^2}}{1+b^2+2b\rho}
\le6\sqrt{\frac{1-\rho}{1+\rho}}.
\]
For $b\le13/20$, the first bound is less than $11/2$.
For $\rho\ge1/11$, the second is at most $\sqrt{30}<11/2$.
Lemma~\ref{lem-square-model-bounds} gives the strict comparison with $8s$.
\end{proof}


\section{Diagonal interpolation and boundary covariance}
\label{sec-global-interpolation}

Use the coordinates of Lemma~\ref{lem-normal-form}, and put
\(d=\sqrt{1-\gamma^2}\). Throughout this section, assume
\begin{equation}
\label{eq-new-compact-affine}
\frac{13}{20}\le b\le1,\qquad
|\gamma|\le\rho:=\frac1{11}.
\end{equation}
Write
\[
A=[e\ f],\quad p=(h,k)^{\mathsf T},\quad
W=A^{\mathsf T}A,\quad
\mathsf G=\det A\,A^{-1}A^{-\mathsf T}
=\frac1d\begin{pmatrix}b&-\gamma\\-\gamma&b^{-1}\end{pmatrix}.
\]
The square functions \(u,v\), the constants \(r,s\), and the boundary
moments \(\ell_0,\ell_1\) are as in
Lemma~\ref{lem-square-model-bounds}.

\subsection{A two-triangle parametrization}

On \(\Omega=\Omega_\square=(-1,1)^2\), put
\[
q(x,y)=1-|x-y|,\qquad n=(1,-1)^{\mathsf T},\qquad
\Phi(x,y)=A\bigl((x,y)^{\mathsf T}+p\,q(x,y)\bigr).
\]
On the two triangles separated by \(x=y\),
\(\nabla q=\pm n\), and
\[
\det D\Phi=\det A\,[1\pm(h-k)]>0.
\]
The map takes these triangles to the two triangles obtained by
cutting \(Q\) along the diagonal joining \(-e-f+g\) and \(e+f+g\).
It is therefore a bi-Lipschitz parametrization of \(Q\).
On \(\partial\Omega\), \(q=xy\).

We transport reference functions to \(Q\) by \(\Phi^{-1}\)
and to \(\Pi=A\Omega\) by \(A^{-1}\).
We use this convention in \eqref{eq-rayleigh-deficit}.

\begin{lemma}[Energy completion on two triangles]
\label{lem-new-tent-energy}
If \(w\in H^1(\Omega)\) is centrally odd, then, for every \(c\in\mathbb R\),
\[
\mathcal E_Q(w+cq)=\mathcal E_\Pi(w)+
\frac{n^{\mathsf T}\mathsf Gn}{1-(h-k)^2}
\int_\Omega(p\cdot\nabla w-c)^2.
\]
\end{lemma}

\begin{proof}
Write \(\xi=\nabla q\) and \(\theta=p\cdot\xi\).
The pulled-back energy matrix is
\[
B_p=(1+\theta)\mathsf G-p\xi^{\mathsf T}\mathsf G
-\mathsf G\xi p^{\mathsf T}
+\frac{\xi^{\mathsf T}\mathsf G\xi}{1+\theta}\,pp^{\mathsf T}.
\]
Under central reflection, \(\nabla w\) is even and \(\xi\) is odd.
Moreover, \(\xi^{\mathsf T}\mathsf G\xi=n^{\mathsf T}\mathsf Gn\)
and \(\theta^2=(h-k)^2\).
Pairing opposite points gives
\[
\mathcal E_Q(w)=\mathcal E_\Pi(w)+
\frac{n^{\mathsf T}\mathsf Gn}{1-(h-k)^2}
\int_\Omega(p\cdot\nabla w)^2.
\]
The identity
\[
B_p\xi=\mathsf G\xi-
\frac{\xi^{\mathsf T}\mathsf G\xi}{1+\theta}p
\]
gives, by the same pairing,
\[
\mathcal E_Q(w,q)=-
\frac{n^{\mathsf T}\mathsf Gn}{1-(h-k)^2}
\int_\Omega p\cdot\nabla w,\qquad
\mathcal E_Q(q)=
\frac{4n^{\mathsf T}\mathsf Gn}{1-(h-k)^2}.
\]
The identity follows by expanding the quadratic energy.
\end{proof}

Set
\[
m=(h-k)^2,\qquad \lambda=1-m,\qquad
R=|g|^2=p^{\mathsf T}Wp,\qquad \delta=1-b.
\]
Thus \(0\le m<1\), \(0<\lambda\le1\), and \(0\le R<1\).
Define
\[
w_\lambda=\lambda u+(1-\lambda)x,\qquad
z_\lambda=\lambda v+(1-\lambda)y,
\]
\[
F_\lambda=w_\lambda+hq,\qquad
G_\lambda=z_\lambda+kq.
\]
We use the two physical trial functions
\[
\mathbf Z\circ\Phi=\begin{pmatrix}Z_1\circ\Phi\\ Z_2\circ\Phi\end{pmatrix}
=A\begin{pmatrix}F_\lambda\\ G_\lambda\end{pmatrix},
\qquad
\mathcal N=D_Q(Z_1)+D_Q(Z_2).
\]
Writing $U=(u,v)^{\mathsf T}$ and $\xi=(x,y)^{\mathsf T}$, we have
\[
\mathbf Z(\Phi(\xi))=\Phi(\xi)+\lambda A\bigl(U(\xi)-\xi\bigr).
\]
The eigenfunction correction is therefore weighted by the product of the
two normalized Jacobian determinants.
The choice of \(\lambda\) cancels the denominator in the energy
completion:
\[
\frac{\lambda^2 n^{\mathsf T}\mathsf Gn}{1-(h-k)^2}
=\frac{\lambda(b+b^{-1}+2\gamma)}d.
\]

\subsection{Boundary covariance and energy}

Set
\[
a_\lambda=\int_{-1}^1w_\lambda(t,1)^2\,dt,\qquad
b_\lambda=\int_{-1}^1w_\lambda(1,t)^2\,dt,
\]
\[
E_\lambda=\int_\Omega|\nabla w_\lambda|^2,\qquad
C_\lambda=\lambda^2C+4\lambda m\ell_0+4m^2.
\]
Then
\[
E_\lambda=\lambda^2E+4\lambda m\ell_0+4m^2,\qquad
E_\lambda-C_\lambda=\lambda^2D_0.
\]

Let
\[
H_1=|e+g|+|e-g|,\quad H_2=|f+g|+|f-g|,
\]
\[
d_1=|e+g|-|e-g|,\quad d_2=|f+g|-|f-g|,
\quad P=2(H_1+H_2).
\]
Put \(\mathbf d=(d_1,d_2)^{\mathsf T}\) and
\(\widehat{\mathbf d}=(d_2,d_1)^{\mathsf T}\), and define
\[
L_0=\lambda\ell_0+2m,\qquad
L_1=\lambda\ell_1+\frac23m.
\]
The boundary identities are
\[
\int_{\partial Q}
\begin{pmatrix}w_\lambda\\z_\lambda\end{pmatrix}\,ds
=L_0\widehat{\mathbf d},\qquad
\int_{\partial Q}
q\begin{pmatrix}w_\lambda\\z_\lambda\end{pmatrix}\,ds
=L_1\mathbf d,
\]
\[
\int_{\partial Q}q\,ds=0,\qquad
\int_{\partial Q}q^2\,ds=\frac P3.
\]
The matrix of uncentered second boundary moments is diagonal, with entries
\[
H_1a_\lambda+H_2b_\lambda,\qquad
H_1b_\lambda+H_2a_\lambda.
\]
These formulas follow by integrating along each reference side,
where \(q=xy\).

Define the displacement contribution
\begin{equation}
\label{eq-new-boundary-functional}
\mathcal B_\lambda=
\frac R3+\frac{2L_1}{P}p^{\mathsf T}W\mathbf d
-\frac{L_0^2}{P^2}
\widehat{\mathbf d}^{\mathsf T}W\widehat{\mathbf d}.
\end{equation}
Let \(\mathcal N_{\rm aff}\) be the normalized deficit of
\(A(w_\lambda,z_\lambda)^{\mathsf T}\) on \(\Pi=A\Omega\).
The change in its normalized uncentered second boundary moment is
\begin{equation}
\label{eq-new-reweight-exact}
\mathcal R_\lambda=
-\frac{\delta(b_\lambda-a_\lambda)}P
\bigl[b(H_1-2)-(H_2-2b)\bigr].
\end{equation}
The energy completion gives
\begin{equation}
\label{eq-new-global-exact}
\mathcal N=
\mathcal N_{\rm aff}+\mathcal R_\lambda+\mathcal B_\lambda
-\frac{\lambda(b+b^{-1}+2\gamma)}{8sd}\,E_{\rm corr},
\end{equation}
where
\begin{equation}
\label{eq-new-Q}
E_{\rm corr}=C_*R-B(1-b^2)(h^2-k^2)-2b\gamma D_0\,hk.
\end{equation}
The correction vector is $(DU-I)p$, so
\[
E_{\rm corr}=\int_\Omega|A(DU-I)p|^2\ge0.
\]
Parity eliminates the odd mixed moments; the remaining moments are
\[
\int_\Omega(\partial_xu-1)^2=C_*-B,\qquad
\int_\Omega(\partial_yu)^2=B,
\]
\[
\int_\Omega\bigl[(\partial_xu-1)(\partial_yv-1)
+\partial_yu\,\partial_xv\bigr]=C_*-D_0.
\]
Substituting these moments gives \eqref{eq-new-Q}.


\subsection{The boundary covariance gain}

The boundary contribution combines the change in side weights with the
variance produced by the displacement.
\begin{lemma}[Boundary covariance]\label{lem-new-combined-boundary}
Under \eqref{eq-new-compact-affine},
\[
\mathcal R_\lambda+\mathcal B_\lambda
\ge\left(\frac9{25}\lambda+\frac3{20}m\right)R.
\]
\end{lemma}
\begin{proof}
The side-length bounds give
\[
H_1-2\le2\sqrt{1+R}-2\le R,\qquad
H_2-2b\ge0,\qquad P\ge4(1+b).
\]
Moreover,
\[
b_\lambda-a_\lambda
=\lambda^2(b_\square-a_\square)
+2\lambda m(\ell_0-\ell_1)+\frac43m^2
\le\frac\lambda3+\frac{4m}{3}.
\]
Here $0<b_\square-a_\square<1/3$ and
$0<\ell_0-\ell_1<5/6$ follow from the square moment bounds.
Since $b(1-b)/(1+b)\le7/50$, \eqref{eq-new-reweight-exact} implies
\begin{equation}\label{eq-boundary-reweight-loss}
\mathcal R_\lambda\ge-\frac7{600}(4-3\lambda)R.
\end{equation}

For $0\le\theta\le1$, set
\[
L_0(\theta)=\theta\ell_0+2(1-\theta),\qquad
L_1(\theta)=\theta\ell_1+\frac23(1-\theta),
\]
\[
V_\theta=\frac R3+\frac{2L_1(\theta)}P
p^{\mathsf T}W\mathbf d
-\frac{L_0(\theta)^2}{P^2}
\widehat{\mathbf d}^{\mathsf T}W\widehat{\mathbf d}.
\]
Thus $V_\lambda=\mathcal B_\lambda$. For fixed geometry,
$V_\theta-7(4-3\theta)R/600$ is concave in $\theta$:
the first moments are affine, and the subtracted square has a
nonnegative coefficient. We prove the endpoint bounds
\begin{equation}\label{eq-boundary-endpoints}
V_0\ge\frac15R,\qquad V_1\ge\frac38R.
\end{equation}

Put
\[
H=H_1+H_2,\qquad a=e\cdot g,\qquad z=\frac fb\cdot g,
\qquad x_1=\frac2{H_1H},\qquad x_2=\frac{2b}{H_2H}.
\]
The two unit vectors $e,f/b$ have scalar product $\gamma$, so
\begin{equation}\label{eq-boundary-frame-comparison}
(1-|\gamma|)R\le a^2+z^2\le(1+|\gamma|)R,
\qquad z^2\le R.
\end{equation}
The identities $H_1d_1=4a$ and $H_2d_2=4bz$ yield
\begin{equation}\label{eq-boundary-displacement-quadratic}
\begin{split}
V_\theta={}&\frac R3+
\bigl(2L_1(\theta)x_1-L_0(\theta)^2b^2x_1^2\bigr)a^2\\
&+\bigl(2L_1(\theta)bx_2-L_0(\theta)^2x_2^2\bigr)z^2
-2\gamma L_0(\theta)^2bx_1x_2az.
\end{split}
\end{equation}
Since $R\le(h+bk)^2<1$,
\[
2\le H_1\le2\sqrt{1+R},\qquad
2b\le H_2\le2\sqrt{b^2+R}.
\]
Consequently,
\begin{equation}\label{eq-boundary-ratio-ranges}
\frac18\le x_1\le\frac1{2(1+b)}<\frac13,\qquad
\frac b{5+4b^2}\le x_2\le\frac1{2(1+b)},\qquad
bx_1x_2\le\frac1{16}.
\end{equation}
For the lower bound on $x_2$, use
\[
H_2H\le4\bigl(1+b^2+\sqrt{2(1+b^2)}\bigr)
\le2(5+4b^2).
\]
The product bound follows from $b/[4(1+b)^2]\le1/16$.

At $\theta=1$, the coefficient of $a^2$ in
\eqref{eq-boundary-displacement-quadratic} exceeds
\[
\frac98x_1-\frac85x_1^2
\ge\frac{71}{120}x_1\ge\frac{71}{960}>\frac1{18}.
\]
The coefficient of $z^2$ is bounded below by the concave quadratic
$f_b(t)=(9/8)bt-(8/5)t^2$. Its endpoint values on the interval
in \eqref{eq-boundary-ratio-ranges} are
\[
\frac{9b}{16(1+b)}-\frac2{5(1+b)^2},\qquad
\frac{b^2\{(9/8)(5+4b^2)-8/5\}}{(5+4b^2)^2}.
\]
Both increase with $b$ and exceed $1/18$ at $b=13/20$.
Using $\ell_0^2bx_1x_2<1/10$ and $2|az|\le a^2+z^2$ gives
\[
V_1\ge\frac R3+
\left(\frac1{18}-\frac1{110}\right)(a^2+z^2)
\ge\frac38R.
\]

At $\theta=0$, the coefficient of $a^2$ is nonnegative. Concavity in
$x_2$ gives
\[
\frac43bx_2-4x_2^2
\ge\min\left\{0,\frac{2b}{3(1+b)}-\frac1{(1+b)^2}\right\}
\ge-\frac{38}{363}.
\]
The last endpoint value increases with $b$.
Since $4bx_1x_2\le1/4$, the two bounds in
\eqref{eq-boundary-frame-comparison} imply
\[
V_0\ge\frac R3-\frac{38}{363}z^2-\frac1{44}(a^2+z^2)
\ge\left(\frac13-\frac{38}{363}-\frac3{121}\right)R
\ge\frac15R.
\]
This proves \eqref{eq-boundary-endpoints}. The concave lower bound
from \eqref{eq-boundary-reweight-loss} therefore yields
\[
\begin{split}
\mathcal R_\lambda+\mathcal B_\lambda
&\ge\left[
\lambda\left(\frac38-\frac7{600}\right)
+m\left(\frac15-\frac7{150}\right)\right]R\\
&\ge\left(\frac9{25}\lambda+\frac3{20}m\right)R.
\end{split}
\]
\end{proof}

\section{The convex deficit}

\begin{lemma}[Affine deficit]
\label{lem-new-affine-baseline}
Under \eqref{eq-new-compact-affine},
\[
\mathcal N_{\rm aff}\ge
-b\left(\frac1s-\frac43\right)m^2+
\delta^2\left(\frac{\lambda^2}{8}
+\frac9{16}\lambda m+\frac{m^2}{3}\right)
+\frac{\gamma^2}{10}.
\]
\end{lemma}

\begin{proof}
First suppose \(\gamma=0\). The normalized weighted second boundary
moment on the parallelogram is
\[
\frac b2(a_\lambda+b_\lambda)
+\frac{\delta^2}{2}a_\lambda.
\]
The weighted energy is
\[
2bE_\lambda+b(b-b^{-1})^2\lambda^2B.
\]
The square eigenfunction identity
\(\mathcal M_\Omega(u,x)=s^{-1}\mathcal E_\Omega(u,x)\)
therefore gives
\[
\mathcal N_{\rm aff}(b,0)=
-b\left(\frac1s-\frac43\right)m^2+
\delta^2\left[
\frac{a_\lambda}{2}
-\frac{\lambda^2B(1+b)^2}{8sb}\right].
\]
Moreover,
\[
a_\lambda=\lambda^2a_\square+2\lambda m\ell_1+\frac23m^2,
\qquad a_\square=\int_{-1}^1u(t,1)^2\,dt>\frac{49}{100}.
\]
Since
\[
\frac{(1+b)^2}{b}=b+2+b^{-1}<\frac{21}{5},
\]
the coefficient of \(\lambda^2\) satisfies
\[
\frac{a_\square}{2}-\frac{B(1+b)^2}{8sb}
>\frac{49}{200}-\frac{63}{550}
=\frac18+\frac3{550}>\frac18.
\]
Together with \(\ell_1>9/16\), this gives the estimate for \(\gamma=0\).

The normalized weighted second boundary moment of the parallelogram is
independent of \(\gamma\). Its weighted energy is
\[
E_{\rm aff}(b,\gamma)=
\frac bd\left[
2E_\lambda+(b-b^{-1})^2\lambda^2B
-2\gamma^2C_\lambda\right].
\]
Consequently,
\[
\begin{split}
E_{\rm aff}(b,0)-E_{\rm aff}(b,\gamma)
={}&\frac{b(1-d)}d
\left[2dC_\lambda
-\lambda^2\bigl(2D_0+(b-b^{-1})^2B\bigr)\right].
\end{split}
\]
The model estimates imply \(C_\lambda>8/5\):
\[
C_\lambda>
\frac85\lambda^2+5\lambda m+4m^2
\ge\frac85(\lambda+m)^2=\frac85.
\]
Also \(d>99/100\) and \((b-b^{-1})^2<21/25\).
Thus the bracket is larger than
\[
2\frac{99}{100}\frac85-\frac14-\frac{21}{25}\frac3{20}
=\frac{349}{125}>2.
\]
The resulting increase in the normalized deficit is at least
\[
\frac{b(1-d)}{4sd}
=\frac{b\gamma^2}{4sd(1+d)}
\ge\frac{b\gamma^2}{8s}
\ge\frac{\gamma^2}{10}.
\]
Here \(b/(8s)>1/10\).
\end{proof}

\begin{theorem}[Convex deficit]\label{thm-uniform-convex}
For every quadrilateral satisfying \eqref{eq-new-compact-affine},
the physical trial pair satisfies
\begin{equation}\label{eq-new-global-positive}
\mathcal N\ge\frac1{200}
\left(\delta^2\lambda^2+\gamma^2+\lambda R+b^2m^2\right).
\end{equation}
In particular, \(\mathcal J(Q)<8s\) unless \(Q\) is a square.
\end{theorem}

\begin{proof}
We first compare the displacement with the mixed term in
\eqref{eq-new-Q}. Cauchy--Schwarz in the metric $W$ gives
\[
\begin{split}
(1+b)^2(h+k)^2
&\le
\frac{(b+b^{-1}+2)(b+b^{-1}+2\rho)}{1-\rho^2}R\\
&\le\frac{121}{12}R\le\frac{256}{25}R.
\end{split}
\]
Here $b+b^{-1}+2<21/5$ and $b+b^{-1}+2\rho<19/8$ on
\eqref{eq-new-compact-affine}. Also
$R\ge2b(1-\rho)hk$, so \eqref{eq-new-Q} and the square moments imply
\[
\begin{split}
E_{\rm corr}
&\le\frac{61}{80}R+
\frac3{20}\delta\sqrt m\,(1+b)(h+k)\\
&\le\frac{61}{80}R+\frac{12}{25}\delta\sqrt{mR}.
\end{split}
\]
The energy coefficient satisfies
\[
\frac{b+b^{-1}+2\gamma}{8sd}
<\frac{19/8}{(11/2)(99/100)}<\frac7{16}.
\]
Applying Lemmas~\ref{lem-new-combined-boundary} and
\ref{lem-new-affine-baseline} in \eqref{eq-new-global-exact}, with
$1/s-4/3<4/33$, gives
\[
\begin{split}
\mathcal N\ge{}&
\frac{\delta^2\lambda^2}{8}+\frac{\gamma^2}{10}
+\left(\frac9{25}-\frac{61}{80}\frac7{16}\right)\lambda R\\
&+\frac9{16}\delta^2\lambda m
-\frac{21}{100}\lambda\delta\sqrt{mR}
+\left(\frac{\delta^2}{3}-\frac{4b}{33}\right)m^2
+\frac3{20}mR.
\end{split}
\]
The mixed terms have the square completion
\[
\frac9{16}\delta^2m-\frac{21}{100}\delta\sqrt{mR}
=\left(\frac34\delta\sqrt m-\frac7{50}\sqrt R\right)^2
-\frac{49}{2500}R,
\]
with residual coefficient
\[
\frac9{25}-\frac{61}{80}\frac7{16}-\frac{49}{2500}
=\left(\frac{33}{400}\right)^2>\frac1{200}.
\]
For the terms containing $m^2$, use
\[
R-b^2m=(1-b^2)h^2+2b(b+\gamma)hk\ge0
\]
and
\[
\frac{(1-b)^2}{3}-\frac{4b}{33}+\frac{3b^2}{20}
\ge\frac13\left(1-\frac{13b}{11}\right)^2+\frac{b^2}{60}.
\]
Consequently,
\[
\begin{split}
\mathcal N\ge{}&
\frac{\delta^2\lambda^2}{8}+\frac{\gamma^2}{10}
+\left(\frac{33}{400}\right)^2\lambda R+\frac{b^2m^2}{60}\\
&+\lambda\left(\frac34\delta\sqrt m-\frac7{50}\sqrt R\right)^2
+\frac{m^2}{3}\left(1-\frac{13b}{11}\right)^2,
\end{split}
\]
which proves \eqref{eq-new-global-positive}.

Both physical traces are nonconstant: their centrally odd parts on the
reference boundary are nonzero linear combinations of the independent
functions $w_\lambda,z_\lambda$. Summing \eqref{eq-rayleigh-deficit}
for the two trials gives $\mathcal J(Q)<8s$ when $\mathcal N>0$.
Since $\lambda>0$, the right side of \eqref{eq-new-global-positive}
vanishes only if $\delta=\gamma=0$ and $R=0$. These conditions give
$h=k=0$, so $Q$ is a square.
\end{proof}

\section{Quantitative stability}

\begin{proof}[Proof of Theorem~\ref{thm-stability}]
Put $S=h^2+k^2$ and $\varepsilon=\sqrt{\delta^2+\gamma^2}$.
The convex deficit controls the quadratic displacement throughout
\eqref{eq-new-compact-affine}, since
\[
R\ge(b^2-b\rho)S\ge\frac13S.
\]
If $\lambda\ge1/2$, the first three terms in
\eqref{eq-new-global-positive} control $S+\varepsilon^2$.
If $\lambda<1/2$, then $m>1/2$; the last term is uniformly positive,
while $S$ and $\varepsilon$ are bounded. On the reference boundary,
$|w_\lambda|,|z_\lambda|\le1$, and
$(1+h)^2+(1+k)^2<5$. Hence
\[
|A(F_\lambda,G_\lambda)^{\mathsf T}|^2
\le(1+\rho)(F_\lambda^2+G_\lambda^2)<6.
\]
The sum of the two normalized boundary variances is therefore at most
$6$, and \eqref{eq-rayleigh-deficit} gives
\begin{equation}\label{eq-new-quadratic-global}
8s-\mathcal J(Q)\ge c_0(S+\varepsilon^2)
\end{equation}
for a fixed $c_0>0$.

The linear affine contribution is detected by the deficit form on
$\operatorname{span}\{u,v\}$. For the transported square functions, write
\[
\mathbf K_Q=
\begin{pmatrix}
D_Q(u,u)&D_Q(u,v)\\
D_Q(u,v)&D_Q(v,v)
\end{pmatrix},
\]
and define $\mathbf K_\Pi$ in the same way on $\Pi=A\Omega$.
Let $M_\square=4/(1+r^2)$. Separated integration gives
\[
\begin{split}
D_\Pi(u,u)&=\frac{a_\square+bb_\square}{2(1+b)}
-\frac{b(E-B)+b^{-1}B}{8sd},\\
D_\Pi(v,v)&=\frac{b_\square+ba_\square}{2(1+b)}
-\frac{bB+b^{-1}(E-B)}{8sd},\\
D_\Pi(u,v)&=\frac{\gamma C}{8sd}.
\end{split}
\]
The square identities
\[
a_\square+b_\square=\frac{M_\square}{2},\qquad
b_\square-a_\square=
\frac{M_\square}{2}\left(\frac r\tau-r^2\right),
\]
\[
E=sM_\square,\qquad
E-2B=\frac{M_\square}{2}
\left[\tau^2(1-r^4)+r^2\right],\qquad
C=\frac{M_\square}{2}(1+r^4)
\]
therefore yield
\begin{equation}\label{eq-deficit-matrix-expansion}
\mathbf K_\Pi=
\begin{pmatrix}
\alpha_0\delta&\beta_0\gamma\\
\beta_0\gamma&-\alpha_0\delta
\end{pmatrix}+O(\varepsilon^2),
\end{equation}
where
\[
\alpha_0=\frac{M_\square}{16}
\left[\frac\tau r(1-r^4)+r^2\right]>0,\qquad
\beta_0=\frac{M_\square(1+r^4)}{16s}>0.
\]

The displacement changes this matrix by $O(S)$. Indeed, the paired
side-length functions satisfy, uniformly near the square,
\[
H_i-2|Ae_i|=O(S),\qquad \mathbf d=O(\sqrt S),\qquad
P(Q)-P(\Pi)=O(S),
\]
where $e_1,e_2$ are the standard coordinate vectors.
The boundary second moments change by $O(S)$, and their centering
correction is quadratic in $\mathbf d$. Lemma~\ref{lem-new-tent-energy},
with $c=0$ and polarization, gives an $O(S)$ change in the energy form.
Thus $\mathbf K_Q-\mathbf K_\Pi=O(S)$ in matrix norm.
The largest eigenvalue of the first matrix in
\eqref{eq-deficit-matrix-expansion} is
$\sqrt{\alpha_0^2\delta^2+\beta_0^2\gamma^2}$. It follows that
\[
\lambda_{\max}(\mathbf K_Q)
=\sqrt{\alpha_0^2\delta^2+\beta_0^2\gamma^2}
+O(\varepsilon^2+S)
\ge c_1\varepsilon-C_1S
\]
in a fixed neighborhood of the square.

Choose a unit eigenvector $\zeta=(\zeta_1,\zeta_2)$ for
$\lambda_{\max}(\mathbf K_Q)$ and put $w=\zeta_1u+\zeta_2v$.
Its trace is nonconstant, and its normalized boundary variance is at most
$2$, since $|w|^2\le u^2+v^2\le2$.
The nonnegativity of $8s-\mathcal J(Q)$ and
\eqref{eq-rayleigh-deficit} imply
\[
8s-\mathcal J(Q)\ge4s\,\lambda_{\max}(\mathbf K_Q)
\ge c_2\varepsilon-C_2S.
\]
Adding $1+C_2/c_0$ times \eqref{eq-new-quadratic-global} to this inequality gives
\[
8s-\mathcal J(Q)\ge c_3(S+\varepsilon)
\]
near the square. Outside that neighborhood, $S+\varepsilon^2$ controls
$S+\varepsilon$ by a fixed positive factor on
\eqref{eq-new-compact-affine}. Combining this with the uniform gap from
Lemma~\ref{lem-diagonal-exclusion} proves the theorem.
\end{proof}

\section{Concave quadrilaterals and completion}

\begin{lemma}\label{lem-concave}
A concave quadrilateral satisfies $\mathcal J(Q)<3\sqrt3$.
A triangle represented with one redundant vertex satisfies
$\mathcal J(Q)\le3\sqrt3$.
\end{lemma}
\begin{proof}
Let the neighbors of the reflex vertex be $a,b$, and let $w$ be the fourth
vertex. The reflex vertex lies inside the triangle $abw$. Normalize
$a=(-1,0)$, $b=(1,0)$, and denote the outer and inner intermediate vertices
by $(x,H)$ and $(z,h)$, where $H>h\ge0$. Let the two broken-line lengths
be $2r_1$ and $2r_2$. Convexity of the sum of distances to $a,b$ gives
$1\le r_2\le r_1$. The outer point lies on an ellipse, so
\[
|Q|=H-h<\sqrt{r_1^2-1}.
\]
For a redundant vertex the inner chain is the segment $ab$, and the last
inequality becomes non-strict, with $r_2=1$ and $h=0$.

The boundary probability measure is the mixture of the two broken-line
measures with weights $r_1/(r_1+r_2)$ and $r_2/(r_1+r_2)$.
Lemma~\ref{lem-broken-line} and the variance-mixture identity show that the
sum $V$ of the two coordinate variances satisfies
\[
V\ge a_*+b_*|Q|^2,\qquad
a_*=\frac{r_1^2-r_1r_2+r_2^2+3}{12},\qquad
b_*=\frac{r_1r_2}{4(r_1+r_2)^2}.
\]
The last term comes from the squared difference of the vertical means,
which is $|Q|^2/4$. Adding the two coordinate Rayleigh inequalities gives
\[
\mathcal J(Q)\le\frac{2A}{a_*+b_*|Q|^2}.
\]
This ratio is increasing for $0\le |Q|\le\sqrt{r_1^2-1}$, since
\[
12(r_1+r_2)^2\bigl(a_*-b_*(r_1^2-1)\bigr)
=(r_1^2-r_1r_2-r_2^2)^2+r_1r_2^2(r_1-r_2)
+3r_1^2+9r_1r_2+3r_2^2>0.
\]
Set $F(r_1,r_2)=a_*+b_*(r_1^2-1)$. With $\eta=r_1-r_2\ge0$,
\[
\partial_{r_2}F=
\frac{2d^4+4r_2d^3+3(r_2^2-1)\eta^2+(7r_2^3-3r_2)\eta+8r_2^4}
{12(r_1+r_2)^3}>0.
\]
Consequently
\[
F(r_1,r_2)\ge F(r_1,1)=\frac{r_1^3+3r_1^2+4}{12(r_1+1)}.
\]
The final comparison factors as
\[
F(r_1,1)^2-\frac4{27}(r_1^2-1)
=\frac{(r_1-2)^2(r_1+2)(r_1+7)(3r_1^2+3r_1+2)}
{432(r_1+1)^2}\ge0.
\]
These inequalities prove the assertion. Strictness in the concave case follows from $|Q|<\sqrt{r_1^2-1}$.
\end{proof}

\begin{proof}[Proof of Theorem~\ref{thm-main}]
For a convex quadrilateral with four genuine corners, apply
Lemma~\ref{lem-diagonal-exclusion} or Theorem~\ref{thm-uniform-convex}. The latter is strict unless $b=1$ and $\gamma=h=k=0$, and these
parameters describe a square. A square attains $8s$ by
Lemma~\ref{lem-square-model-bounds}. Every remaining simple quadrilateral
is concave or a triangle with a redundant vertex; apply
Lemma~\ref{lem-concave}, noting that $3\sqrt3<11/2<8s$.
The fixed-perimeter statement follows from similarity invariance.
\end{proof}

\appendix
\section{Elementary square moments}\label{app-square}

The root relation gives \(r=\cot\tau\). Separated integration yields
\begin{align*}
a_\square&=1-\frac{\sin(2\tau)}{2\tau},
&a_\square+b_\square&=\frac2{1+r^2},\\
\ell_0&=\frac{2\cos\tau}{\tau},
&\ell_1&=\frac{2\sin\tau(1-s)}{\tau^2},\\
E&=2\tau\sin(2\tau),
&C&=\frac{2(1+r^4)}{1+r^2},\\
B&=\frac E2-\tau^2(1-r^2)-\frac{r^2}{1+r^2},
&\int_{\Omega_\square}\partial_xu&=2\ell_0.
\end{align*}

Put
\[
f(t)=\cos t\cosh t-\sin t\sinh t,\qquad
p(t)=1-t^2-\frac{t^4}{6}+\frac{t^6}{90}+\frac{t^8}{2520}.
\]
The alternating series for the two products give, on \([0,1]\),
\[
p(t)-\frac{t^{10}}{113400}-\frac{t^{12}}{7484400}
\le f(t)\le p(t).
\]
For \(\tau_-=15/16\) and \(\tau_+=586/625\),
\[
p(\tau_-)-\frac{\tau_-^{10}}{113400}
-\frac{\tau_-^{12}}{7484400}>0>p(\tau_+).
\]
Since \(f'(t)=-2\sin t\cosh t<0\), we have
\(\tau_-<\tau<\tau_+\).

Define the Taylor polynomials
\[
S_j(t)=\sum_{i=0}^j\frac{(-1)^it^{2i+1}}{(2i+1)!},\qquad
C_j(t)=\sum_{i=0}^j\frac{(-1)^it^{2i}}{(2i)!}.
\]
On \([0,2]\),
\[
S_3(t)\le\sin t\le S_4(t),\qquad
C_3(t)\le\cos t\le C_4(t).
\]
Since $\cot t$ decreases on $[\tau_-,\tau_+]$,
\[
\frac{C_3(\tau_+)}{S_4(\tau_+)}<r<
\frac{C_4(\tau_-)}{S_3(\tau_-)},\qquad
\frac7{13}<r^2<\frac{27}{50},\qquad
\frac{11}{16}<s<\frac{69}{100}.
\]
The remaining estimates follow from these parameter bounds. In fact,
$r/(1+r^2)$ increases for $0<r<1$, so
\[
a_\square=1-\frac{r}{\tau(1+r^2)}
>1-\frac{\sqrt{27/50}}{\tau_-(1+27/50)}>\frac{49}{100},
\]
\[
E=\frac{4\tau r}{1+r^2}
<\frac{4\tau_+\sqrt{27/50}}{1+27/50}<\frac{43}{24}.
\]
Also $\ell_0^2=4r^2/[\tau^2(1+r^2)]$, whence
\[
\left(\frac{121}{96}\right)^2
<\frac7{5\tau_+^2}<\ell_0^2
<\frac{108}{77\tau_-^2}<\frac85.
\]
The identity
\[
b_\square-a_\square
=\frac{2r}{1+r^2}\left(\frac1\tau-r\right)>0
\]
and \(a_\square+b_\square<13/10\) give
\(a_\square\le b_\square<81/100\).
Green's identity for $u$ and the coordinate function $x$ gives
\[
2\ell_0=2s(\ell_0+\ell_1),\qquad
\ell_1=(s^{-1}-1)\ell_0.
\]
Consequently,
\[
\ell_1>\frac{121}{96}\frac{31}{69}>\frac9{16},\qquad
\ell_1<\frac43\frac5{11}<\frac23.
\]
The bounds on \(E\) and \(\ell_0\) imply
\[
C_*=E-4\ell_0+4<\frac{43}{24}-\frac{121}{24}+4=\frac34.
\]
Also,
\[
0<B<\frac9{10}-\left(\frac{15}{16}\right)^2\frac{23}{50}
-\frac7{20}<\frac3{20}.
\]
Since \(r^2>1/2\),
\[
C-\frac53=
\frac{(3r^2-1)(2r^2-1)}{3(1+r^2)}>0.
\]
Finally,
\[
D_0=\frac{2(2s-1-r^4)}{1+r^2}>0,\qquad
D_0=E-C<\frac{43}{24}-\frac53=\frac18.
\]
The first inequality follows from \(2s>11/8\) and
\(r^4<(27/50)^2<3/8\).

\begin{thebibliography}{99}
\raggedbottom

\bibitem{AkhmetgaliyevEtAl2017}
E. Akhmetgaliyev, C.-Y. Kao, and B. Osting,
\href{https://doi.org/10.1137/16M1067263}{\emph{Computational methods for extremal Steklov problems}},
SIAM J. Control Optim. \textbf{55} (2017), no.~2, 1226--1240.

\bibitem{AuchmutyCho2015}
G. Auchmuty and M. Cho,
\href{https://doi.org/10.1080/01630563.2015.1031383}{\emph{Boundary integrals and approximations of harmonic functions}},
Numer. Funct. Anal. Optim. \textbf{36} (2015), no.~6, 687--703.

\bibitem{Bogosel2016}
B. Bogosel,
\href{https://doi.org/10.1016/j.cam.2016.04.008}{\emph{The method of fundamental solutions applied to boundary eigenvalue problems}},
J. Comput. Appl. Math. \textbf{306} (2016), 265--285.

\bibitem{Bogosel2017}
B. Bogosel,
\href{https://doi.org/10.1007/s00245-015-9321-5}{\emph{The Steklov spectrum on moving domains}},
Appl. Math. Optim. \textbf{75} (2017), no.~1, 1--25.

\bibitem{BogoselBucur2024}
B. Bogosel and D. Bucur,
\href{https://doi.org/10.5802/jep.250}{\emph{On the polygonal Faber--Krahn inequality}},
J. \'Ec. polytech. Math. \textbf{11} (2024), 19--105.

\bibitem{BogoselBucurGiacomini2017}
B. Bogosel, D. Bucur, and A. Giacomini,
\href{https://doi.org/10.1137/16M1075260}{\emph{Optimal shapes maximizing the Steklov eigenvalues}},
SIAM J. Math. Anal. \textbf{49} (2017), no.~2, 1645--1680.

\bibitem{BrascoEtAl2012}
L. Brasco, G. De Philippis, and B. Ruffini,
\href{https://doi.org/10.1016/j.jfa.2012.03.017}{\emph{Spectral optimization for the Stekloff--Laplacian: the stability issue}},
J. Funct. Anal. \textbf{262} (2012), no.~11, 4675--4710.

\bibitem{Brock2001}
F. Brock,
\href{https://doi.org/10.1002/1521-4001(200101)81:1\%3C69::AID-ZAMM69\%3E3.0.CO;2-\%23}{\emph{An isoperimetric inequality for eigenvalues of the Stekloff problem}},
ZAMM Z. Angew. Math. Mech. \textbf{81} (2001), no.~1, 69--71.

\bibitem{BucurEtAl2021}
D. Bucur, V. Ferone, C. Nitsch, and C. Trombetti,
\href{https://doi.org/10.4310/JDG/1620272940}{\emph{Weinstock inequality in higher dimensions}},
J. Differential Geom. \textbf{118} (2021), no.~1, 1--21.

\bibitem{BucurFragala2021}
D. Bucur and I. Fragal\`a,
\href{https://doi.org/10.1051/cocv/2020083}{\emph{Symmetry results for variational energies on convex polygons}},
ESAIM Control Optim. Calc. Var. \textbf{27} (2021), Art.~3.

\bibitem{BucurNahon2021}
D. Bucur and M. Nahon,
\href{https://doi.org/10.1090/tran/8302}{\emph{Stability and instability issues of the Weinstock inequality}},
Trans. Amer. Math. Soc. \textbf{374} (2021), no.~3, 2201--2223.

\bibitem{ChengEtAl2026}
Z. Cheng, C. Gui, Y. Hu, Q. Li, and R. Yao,
\emph{Global strict monotonicity and asymptotics of the first Steklov eigenvalue for regular polygons}, preprint,
\href{https://arxiv.org/abs/2603.25116v2}{arXiv:2603.25116v2} (2026).

\bibitem{ColboisEtAl2024}
B. Colbois, A. Girouard, C. Gordon, and D. Sher,
\href{https://doi.org/10.1007/s13163-023-00480-3}{\emph{Some recent developments on the Steklov eigenvalue problem}},
Rev. Mat. Complut. \textbf{37} (2024), no.~1, 1--161.

\bibitem{DrydenEtAl2025}
E. B. Dryden, C. Gordon, J. Moreno, J. Rowlett, and C. Villegas-Blas,
\href{https://doi.org/10.1007/s12220-025-01922-8}{\emph{The Steklov spectrum of convex polygonal domains I: Spectral finiteness}},
J. Geom. Anal. \textbf{35} (2025), no.~3, Art.~91.

\bibitem{DrydenEtAl2026}
E. B. Dryden, C. Gordon, J. Moreno, J. Rowlett, and C. Villegas-Blas,
\href{https://doi.org/10.1007/s12220-026-02597-5}{\emph{The Steklov spectrum of convex polygonal domains II: Investigating spectral determination}},
J. Geom. Anal. \textbf{36} (2026), no.~11, Art.~346.

\bibitem{FraserSchoen2011}
A. Fraser and R. Schoen,
\href{https://doi.org/10.1016/j.aim.2010.11.007}{\emph{The first Steklov eigenvalue, conformal geometry, and minimal surfaces}},
Adv. Math. \textbf{226} (2011), no.~5, 4011--4030.

\bibitem{FraserSchoen2016}
A. Fraser and R. Schoen,
\href{https://doi.org/10.1007/s00222-015-0604-x}{\emph{Sharp eigenvalue bounds and minimal surfaces in the ball}},
Invent. Math. \textbf{203} (2016), no.~3, 823--890.

\bibitem{GirouardKarpukhinLagace2021}
A. Girouard, M. Karpukhin, and J. Lagac\'e,
\href{https://doi.org/10.1007/s00039-021-00573-5}{\emph{Continuity of eigenvalues and shape optimisation for Laplace and Steklov problems}},
Geom. Funct. Anal. \textbf{31} (2021), no.~3, 513--561.

\bibitem{GirouardEtAl2019}
A. Girouard, J. Lagac\'e, I. Polterovich, and A. Savo,
\href{https://doi.org/10.1112/S0025579318000414}{\emph{The Steklov spectrum of cuboids}},
Mathematika \textbf{65} (2019), no.~2, 272--310.

\bibitem{GirouardPolterovich2010}
A. Girouard and I. Polterovich,
\href{https://doi.org/10.1007/s10688-010-0014-1}{\emph{On the Hersch--Payne--Schiffer inequalities for Steklov eigenvalues}},
Funct. Anal. Appl. \textbf{44} (2010), no.~2, 106--117.

\bibitem{GirouardPolterovich2017}
A. Girouard and I. Polterovich,
\href{https://doi.org/10.4171/JST/164}{\emph{Spectral geometry of the Steklov problem}},
J. Spectr. Theory \textbf{7} (2017), no.~2, 321--359.

\bibitem{GuiHuLi2026}
C. Gui, Y. Hu, and Q. Li,
\emph{A quantitative P\'olya--Szeg\H{o} theorem for tangential polygons}, preprint,
\href{https://arxiv.org/abs/2607.28768v2}{arXiv:2607.28768v2} (2026).

\bibitem{GuiHuLiYang2026}
C. Gui, Y. Hu, Q. Li, and H. Yang,
\emph{A quantitative Robin P\'olya--Szeg\H{o} theorem for tangential polygons}, preprint,
\href{https://arxiv.org/abs/2610.21417}{arXiv:2610.21417} (2026).

\bibitem{GuiHuLiZhang2026}
C. Gui, Y. Hu, Q. Li, and C. Zhang,
\emph{Mixed torsion on right triangles and the P\'olya--Szeg\H{o} monotonicity problem for regular polygons}, preprint,
\href{https://arxiv.org/abs/2606.13448}{arXiv:2606.13448} (2026).

\bibitem{HerschPayneSchiffer1974}
J. Hersch, L. E. Payne, and M. M. Schiffer,
\href{https://doi.org/10.1007/BF00248412}{\emph{Some inequalities for Stekloff eigenvalues}},
Arch. Rational Mech. Anal. \textbf{57} (1974), no.~2, 99--114.

\bibitem{KrymskiEtAl2021}
S. Krymski, M. Levitin, L. Parnovski, I. Polterovich, and D. A. Sher,
\href{https://doi.org/10.1093/imrn/rnaa200}{\emph{Inverse Steklov spectral problem for curvilinear polygons}},
Int. Math. Res. Not. IMRN \textbf{2021} (2021), no.~1, 1--37.

\bibitem{LevitinEtAl2022}
M. Levitin, L. Parnovski, I. Polterovich, and D. A. Sher,
\href{https://doi.org/10.1112/plms.12461}{\emph{Sloshing, Steklov and corners: Asymptotics of Steklov eigenvalues for curvilinear polygons}},
Proc. Lond. Math. Soc. (3) \textbf{125} (2022), no.~3, 359--487.

\bibitem{Polya1948}
G. P\'olya,
\href{https://doi.org/10.1090/qam/26817}{\emph{Torsional rigidity, principal frequency, electrostatic capacity and symmetrization}},
Quart. Appl. Math. \textbf{6} (1948), no.~3, 267--277.

\bibitem{PolyaSchiffer1954}
G. P\'olya and M. Schiffer,
\href{https://doi.org/10.1007/BF02803593}{\emph{Convexity of functionals by transplantation}},
J. Analyse Math. \textbf{3} (1953--1954), 245--345.

\bibitem{Szego1954}
G. Szeg\H{o},
\href{https://doi.org/10.1512/iumj.1954.3.53017}{\emph{Inequalities for certain eigenvalues of a membrane of given area}},
J. Rational Mech. Anal. \textbf{3} (1954), no.~3, 343--356.

\bibitem{Tan2017}
A. Tan,
\emph{The Steklov problem on rectangles and cuboids}, preprint,
\href{https://arxiv.org/abs/1711.00819}{arXiv:1711.00819} (2017).

\bibitem{Weinberger1956}
H. F. Weinberger,
\href{https://doi.org/10.1512/iumj.1956.5.55021}{\emph{An isoperimetric inequality for the N-dimensional free membrane problem}},
J. Rational Mech. Anal. \textbf{5} (1956), no.~4, 633--636.

\bibitem{Weinstock1954}
R. Weinstock,
\href{https://doi.org/10.1512/iumj.1954.3.53036}{\emph{Inequalities for a classical eigenvalue problem}},
J. Rational Mech. Anal. \textbf{3} (1954), no.~6, 745--753.

\end{thebibliography}
\end{document}
